% TOP_PROBLEM: {"id": 4800010, "title": "Multiple ergodic averages — Problem 10", "category_id": 11, "category": {"id": 11, "name": "dynamical_systems", "display_name": "Dynamical Systems", "description": "Problems about long-term behavior of deterministic systems, Hamiltonian dynamics, and periodic orbits.", "slug": "dynamical-systems", "order_index": 11, "created_at": "2026-07-31T15:26:25.671Z"}, "status": "partially_solved", "problem_number": "AMR-047-0010", "set_id": 13}
% FIRST_POSTED: 2026-10-01
% ENTRY_KIND: statement_only
% UnsolvedMath ID 4800010; source catalogue code AMR-047-0010
% !TeX program = lualatex
% Definition and sources only; this file is not an LLM solution attempt.
\documentclass[11pt,a4paper]{article}
\usepackage[margin=25mm]{geometry}
\usepackage{fontspec}
\setmainfont{lmroman10-regular.otf}[BoldFont=lmroman10-bold.otf,
 ItalicFont=lmroman10-italic.otf,BoldItalicFont=lmroman10-bolditalic.otf]
\usepackage{amsmath,amssymb,mathtools}
\usepackage{unicode-math}
\setmathfont{latinmodern-math.otf}
\AtBeginDocument{\let\setminus\smallsetminus}
\usepackage{xurl,hyperref}
\hypersetup{colorlinks=true,urlcolor=blue}
\setcounter{secnumdepth}{0}
\setlength{\parindent}{0pt}
\setlength{\parskip}{5pt}
\setlength{\emergencystretch}{3em}
\begin{document}
\section*{Multiple ergodic averages — Problem 10}\label{4800010}
{\small UnsolvedMath ID 4800010 · AMR-047-0010 · Dynamical Systems · AMR Open Problem Lists}
\subsection{Definitions and mathematical statement}
Fix $\ell\ge1$. For each $N\ge1$ let
$p_{1,N},\ldots,p_{\ell,N}\in\mathbb R[t]$, with degrees
bounded by one finite constant independent of $i,N$.
Call this array good if for every nonzero vector
$c=(c_1,\ldots,c_\ell)\in\mathbb R^\ell$,
\[
       \lim_{N\to\infty}\frac1N\sum_{n=1}^N
           \exp\!\left(2\pi i\sum_{j=1}^\ell
                                   c_jp_{j,N}(n)\right)=0.
\]
This is the source's requirement that every nontrivial
real linear combination be a good polynomial sequence;
the coefficients $c_j$ are fixed as $N$ varies.

Let $T$ be any invertible ergodic measure-preserving
transformation of a probability space $(X,\mathcal B,\mu)$.
Prove that for every $f_1,\ldots,f_\ell\in L^\infty(\mu)$,
\[
 \frac1N\sum_{n=1}^N\prod_{j=1}^\ell
       f_j\circ T^{\lfloor p_{j,N}(n)\rfloor}
       \longrightarrow\prod_{j=1}^\ell\int_X f_j\,d\mu
                        \quad\text{in }L^2(\mu).
\]
The floor denotes the greatest integer at most its real
argument, including for negative values. Ergodicity means
that invariant measurable sets have measure zero or one;
total ergodicity is not assumed. Polynomial coefficients
may vary with $N$, without separate size bounds. The
same upper index $N$ specifies both the polynomial row
and the averaging interval. The conclusion asks for the
displayed product of integrals, not merely existence of
some limit.
\subsection{Short English statement}
Does joint exponential-sum cancellation for variable polynomials force the corresponding ergodic averages to converge to the product of means?
\subsection{Sources}
\begin{itemize}
\item[S1] ULAM.AI and the UnsolvedMath Contributors, supplied problems.json, record 4800010 (AMR-047-0010), AMR Open Problem Lists. Catalogue identity and source transcription; CC BY 4.0.
\url{https://huggingface.co/datasets/ulamai/UnsolvedMath}.
\item[S2] Frantzikinakis - Some open problems on multiple ergodic averages (2016), Problem 10. Original source indexed by AMR-047-0010.
\url{https://arxiv.org/abs/1103.3808}.
\end{itemize}
\end{document}
